% OR03-S012 | lapisan asli: laboratorium 3 globalisasi Newton
% segment-id: d90.orig.v1.tr03.seg0012
% lab-id: d90.orig.v1.tr03.lab.0003
% lab-hint-id: d90.orig.v1.tr03.lab.hint.0003
% lab-answer-id: d90.orig.v1.tr03.lab.answer.0003
% lab-solution-id: d90.orig.v1.tr03.lab.solution.0003

\section{Laboratorium 3: kegagalan lokal dan globalisasi Newton}
\label{orig03:lab:globalisasi-newton}

Laboratorium ini memperlihatkan bahwa arah yang tampak ``orde dua'' belum
tentu aman secara global. Kita memakai fungsi konveks mulus
\begin{equation}
  F(x,y)=\log(\cosh x)+\frac12y^2,
  \qquad
  \nabla F(x,y)=\begin{pmatrix}\tanh x\\y\end{pmatrix},
  \qquad
  \nabla^2F(x,y)=
  \begin{pmatrix}\operatorname{sech}^2x&0\\0&1\end{pmatrix}.
  \label{orig03:eq:lab3-model}
\end{equation}
Minimum uniknya adalah $(0,0)$ dengan nilai nol. Walaupun Hessian selalu
positif, kelengkungan pada arah $x$ sangat kecil ketika $|x|$ besar. Karena
itu, langkah Newton penuh dari titik awal yang disediakan dapat melompat ke
daerah dengan nilai objektif jauh lebih tinggi. Contoh ini memisahkan dua
persoalan: pemilihan arah dan pemilihan panjang langkah.

Kode tetap berada di
\texttt{labs/original-03/globalisasi-newton.py}. Dengan seed dan konfigurasi
yang dibekukan, kode membandingkan empat metode:
\begin{enumerate}
  \item gradien dengan langkah tetap yang sengaja terlalu besar;
  \item gradien dengan pencarian balik Armijo;
  \item Newton murni dengan langkah penuh;
  \item Newton teredam dengan Hessian terkoreksi
    $B_k=\nabla^2F(z_k)+\tau_kI$, dengan
    $\lambda_{\min}(B_k)\geq \underline\mu$.
\end{enumerate}
Jejak lengkap ditulis ke CSV, ringkasan dan sertifikat ditulis ke JSON, dan
SVG hanya memvisualkan data yang sama. Tidak ada kesimpulan numerik yang
bergantung pada grafik.

\subsection*{Prompt}
\begin{exercise}[Tugas laboratorium]
Jalankan skrip sekali dan kerjakan hal-hal berikut.
\begin{enumerate}
  \item Turunkan langkah Newton murni pada koordinat $x$ dan jelaskan mengapa
    $x_0\approx3$ menghasilkan lompatan besar walaupun fungsi konveks.
  \item Dari jejak CSV, temukan kenaikan objektif pertama bagi Newton murni
    dan kegagalan metode gradien berlangkah tetap.
  \item Buktikan bahwa syarat Armijo
    \[
      F(z_k+\alpha_kd_k)\leq F(z_k)+c_1\alpha_k
      \nabla F(z_k)^\top d_k
    \]
    menghasilkan penurunan ketat bila $d_k$ adalah arah turun.
  \item Buktikan bahwa koreksi
    $\tau_k=\max\{0,\underline\mu-\lambda_{\min}(\nabla^2F(z_k))\}$
    membuat $d_k=-B_k^{-1}\nabla F(z_k)$ menjadi arah turun selama
    gradien tidak nol.
  \item Cocokkan bukti Anda dengan seluruh boolean sertifikat dalam JSON dan
    jelaskan mengapa ``gagal'' di sini merupakan bukti terkontrol, bukan
    galat program.
\end{enumerate}
\end{exercise}

\subsection*{Petunjuk bertahap}
\label{orig03:lab:hint:0003}
\begin{enumerate}
  \item Gunakan $\operatorname{sech}^2x=1-\tanh^2x$ dan tulis
    $x^+=x-\tanh(x)/\operatorname{sech}^2(x)$.
  \item Bandingkan kolom \texttt{objective} pada dua baris berturutan; jangan
    menyimpulkan kegagalan hanya dari gambar.
  \item Pada ruas kanan Armijo, tanda
    $\nabla F(z_k)^\top d_k<0$ adalah fakta penentu.
  \item Gunakan $B_k\succeq\underline\mu I$ untuk menunjukkan
    $\nabla F(z_k)^\top d_k=-g_k^\top B_k^{-1}g_k<0$.
  \item Periksa juga bahwa semua langkah metode teredam diterima oleh Armijo
    dan nilai objektifnya monoton menurun.
\end{enumerate}

\subsection*{Jawaban singkat}
\label{orig03:lab:answer:0003}
Newton murni gagal karena invers kelengkungan
$1/\operatorname{sech}^2(x_0)$ sangat besar. Langkah gradien tetap gagal
karena faktor koordinat kuadratiknya adalah $1-\alpha=-1.5$. Pencarian balik
memendekkan langkah sampai memenuhi penurunan cukup. Koreksi Hessian menjamin
arah turun, sedangkan Armijo mengglobalisasi panjang langkah; kombinasi
keduanya mencapai toleransi gradien yang dibekukan.

\subsection*{Solusi acuan lengkap}
\label{orig03:lab:solution:0003}
Pada koordinat pertama, Newton murni memberi
\[
 x^+=x-\frac{\tanh x}{\operatorname{sech}^2x}
     =x-\sinh x\cosh x.
\]
Untuk $x\approx3$, suku $\sinh x\cosh x$ berorde $10^2$, sehingga iterasi
melompat melintasi minimizer dan menaikkan $\log(\cosh x)$. Pada iterasi
berikutnya $\operatorname{sech}^2x$ mendekati nol pada presisi mesin; skrip
berhenti sebelum pembagian hampir singular merusak jejak. Untuk gradien tetap
pada koordinat kedua berlaku $y^+=(1-2.5)y=-1.5y$, jadi magnitudonya tumbuh
geometrik.

Untuk arah gradien $d=-g$, kita memiliki $g^\top d=-\|g\|^2<0$. Karena $F$
diferensiabel, terdapat $\bar\alpha>0$ sehingga setiap
$0<\alpha\leq\bar\alpha$ memenuhi Armijo; pencarian balik dengan faktor dalam
$(0,1)$ karena itu berhenti dalam banyak langkah berhingga. Untuk arah Newton
terkoreksi, konstruksi $B_k\succeq\underline\mu I$ memberi
\[
 g_k^\top d_k=-g_k^\top B_k^{-1}g_k
 \leq-\frac{\|g_k\|^2}{\lambda_{\max}(B_k)}<0.
\]
Maka argumen pencarian balik yang sama berlaku. Sertifikat program memeriksa
langsung batas eigenvalue, tanda turunan arah, semua penerimaan Armijo,
monotonisitas objektif, dan norma gradien akhir. Sebaliknya, metode yang tidak
diglobalisasi diwajibkan memperlihatkan kenaikan atau ledakan yang telah
ditentukan; dengan demikian jalur gagal merupakan bagian dari eksperimen.

\paragraph{Hak dan firewall O018.}
Teks, soal, solusi, dan kode laboratorium ini ditulis mandiri untuk lapisan
Original-03 dan tersedia berdasarkan CC BY-SA 4.0. Tidak ada prosa, soal,
solusi, kode, atau susunan pedagogis O018 yang disalin atau diadaptasi; O018
tetap merupakan korpus terpisah.
